\centerline{\bf \Huge 5 класс}

\textbf{}

\problem{5.1} Коробка и 100 лежащих в ней одинаковых алюминиевых
шариков весят 510 г. Такая же коробка и 100 одинаковых пластмассовых шариков, лежащих в ней, весят 492 г. Сколько будет весить коробка и лежащие в ней 20 алюминиевых и 80 пластмассовых шариков? 

\problem{5.2} Волк, Ёж, Чиж и Бобёр делили апельсин. Ежу досталось вдвое больше долек, чем Чижу, Чижу — впятеро меньше, чем Бобру, а Бобру — на 8 долек больше, чем Чижу. Найдите, сколько долек было в апельсине, если Волку досталась только кожура.

\problem{5.3} Аня лжет по вторникам, средам и четвергам и говорит правду во все остальные дни недели. Ваня же лжёт по четвергам, пятницам и субботам и говорит правду во все остальные дни недели. На вопрос: «Какой сегодня день недели?» Аня ответила: «Пятница», а Ваня – «Вторник». В какие дни недели это могло случиться? Ответ нужно объяснить.

\problem{5.4} Как разрезать прямоугольник длиной в 18 см и шириной 8 см на наименьшее число прямоугольников и сложить из них квадрат?

\newpage

\centerline{\bf \Huge 6 класс}

\textbf{}

\problem{6.1} В двузначном числе A поменяли цифры местами и получили число B. Найдите такое A, чтобы сумма A+B делилась на 17.

\problem{6.2} В каждом из трех сундуков Али-Баба нашел золотые и серебряные монеты; всего 40 золотых и 40 серебряных монет. В первом сундуке золотых монет было на 7 больше, чем серебряных, во втором —серебряных на 15 меньше, чем золотых. Каких монет больше в третьем сундуке и на сколько?

\problem{6.3} Два велосипедиста собрались доехать из пункта A в пункт B. Скорость первого из них равна 35 км/ч, второго – 25 км/ч. Известно, что каждый из них ехал только тогда, когда другой отдыхал (стоял на месте), а всего за 2 часа они проехали одинаковое расстояние. Могли ли они доехать за 2 часа до пункта B, расположенного на расстоянии 30 км от пункта A?

\problem{6.4} За круглый стол сели 12 человек, некоторые из них — рыцари, а остальные — лжецы (рыцари всегда говорят правду, а лжецы всегда лгут). Затем каждый из них сказал: «Среди моих соседей есть лжец». Какое наибольшее число из сидящих за столом может сказать: «Среди моих соседей есть рыцарь»?

\newpage

\centerline{\bf \Huge 7 класс}

\textbf{}

\problem{7.1} Расставьте по кругу 6 ненулевых цифр (не обязательно различных) так, чтобы каждая из них равнялась последней цифре суммы своих соседей. 
  
\problem{7.2} Рационализация процесса производства дала возможность увеличить в октябре выпуск пластмассовых контейнеров на 20\% по сравнению с сентябрем, в ноябре на 5\% по сравнению с октябрем и в декабре на 10\% по сравнению с ноябрем. В результате в декабре цех выпустил 11088 контейнеров. Сколько контейнеров выпустил цех в ноябре, в октябре, в сентябре?

\problem{7.3} Поменяйте в выражении $2025^2 * 2024^2 * \dotsc * 2^2 * 1^2$ звездочки на плюс или минус так, чтобы получилось 2025.

\problem{7.4} Дорожная система Элиста представляет собой 15 бесконечных параллельных проспектов, пересекающих через равные расстояния поперечные улицы. В городе завелся бандит Гольдан, который прячется за домами. За законом в городе следят шериф Расул и маршалл Анджур. Если бандит окажется на одной улице или на одном проспекте с шерифом или маршалом, то он тут же будет пойман. Скорость Гольдана не более чем в 15 раз выше чем скорость Расула и Анджура, причем они знают, что Гольдан в начальный момент времени находился от них на расстоянии не более чем 225 кварталов. Докажите, что защитники закона справятся с заданием и спасут Элисту от бандита.

\newpage

\centerline{\bf \Huge 8 класс}

\textbf{}

\problem{8.1} Найдите минимальное натуральное $n$ при котором $n!$ делится на 715.

\problem{8.2}Герман и Ердем сделали в тире по $5$ выстрелов. Первыми тремя выстрелами они выбили поровну, а последними тремя Герман выбил в три раза больше очков, чем Ердем. На мишени остались пробоины в $10, 9, 9, 8, 8, 5, 4, 4, 3, 2$ очков. Куда попал каждый из них третьим выстрелом? Приведите все возможные варианты ответа и докажите, что других нет


\problem{8.3} В прямоугольнике $ABCD$ точка $M$ — середина стороны $BC$, точка $N$ — середина стороны $CD$, $P$ — точка пересечения отрезков $DM$ и $BN$. Докажите, что угол $MAN$ равен углу $BPM$. 

\problem{8.4} Елена Михайловна вызвала Егора к доске и попросила нарисовать график функции $y = \dfrac{1}{x}.$ Затем она вызвала к доске Глеба и сказала ему отмечать все точки графика у которых на оси ОХ координата 1, 2, 3 и т.д. пока он не устанет. Глеб принялся за дело и отметил ровно n точек. После Глеба к доске вышла Алина и начала закрашивать все прямоугольники у которых одна вершина это точка отмеченная Глебом, противоположная вершина это начало координат, а остальные вершины это лежат на осях(смотреть на рисунок). После всего этого безобразия Елена Михайловна попросила детей посчитать площадь той части, что была покрашена Алиной только один раз. Какой должен получиться ответ?

а) При n = 10

б) При n произвольном.

\begin{figure}[h] 
    \hspace*{-13cm}
    \includegraphics[scale=2]{Осень_2023/pictureolymp.png}
\end{figure}